% Created 2026-06-01 Mon 23:52
% Intended LaTeX compiler: pdflatex
\documentclass[11pt, letterpaper]{article}
\usepackage[utf8]{inputenc}
\usepackage[T1]{fontenc}
\usepackage{graphicx}
\usepackage{longtable}
\usepackage{wrapfig}
\usepackage{rotating}
\usepackage[normalem]{ulem}
\usepackage{amsmath}
\usepackage{amssymb}
\usepackage{capt-of}
\usepackage{hyperref}
\usepackage[margin=2.5cm]{geometry}
\usepackage[utf8]{inputenc}
\usepackage[T1]{fontenc}
\usepackage{lmodern}
\usepackage{xcolor}
\usepackage{hyperref}
\hypersetup{colorlinks=true, linkcolor=blue, urlcolor=blue}
\author{Equipo Alpine White}
\date{\textit{{[}2026-05-25 Mon]}}
\title{Bitácora Semanal XX: Administración de Sistemas Unix/Linux}
\hypersetup{
 pdfauthor={Equipo Alpine White},
 pdftitle={Bitácora Semanal XX: Administración de Sistemas Unix/Linux},
 pdfkeywords={},
 pdfsubject={},
 pdfcreator={Emacs 30.2 (Org mode 9.7.39)}, 
 pdflang={English}}
\begin{document}

\maketitle
\setcounter{tocdepth}{2}
\tableofcontents

\section{Información General}
\label{sec:org857ddcd}

\begin{itemize}
\item \textbf{Semana:} Semana XX
\item \textbf{Materia:} Administración de Sistemas Unix/Linux
\item \textbf{Docente:} Luis Enrique Serrano Gutiérrez
\item \textbf{Equipo:}

\begin{itemize}
\item Arreguín Salgado Gael Emiliano
\item López Pérez Mariana
\item Nieto Gallegos Isaac Julián
\end{itemize}
\end{itemize}

Durante esta semana se trabajó con tecnologías y herramientas fundamentales para la administración avanzada de sistemas Linux, enfocándonos principalmente en comprender cómo el sistema operativo abstrae, organiza y administra recursos críticos como almacenamiento, servicios de red, sincronización, permisos y procesos internos.

A lo largo de las prácticas se observó cómo Linux mantiene una filosofía basada en modularidad y automatización, permitiendo que múltiples tecnologías colaboren entre sí de forma transparente.

Más allá de aprender únicamente comandos, el objetivo principal fue entender la lógica detrás de cada herramienta, su impacto sobre el sistema y las implicaciones administrativas y de seguridad que conlleva su configuración.

La práctica permitió reforzar conceptos relacionados con almacenamiento, sistemas de archivos, servicios de red y administración dinámica de recursos. También se exploró cómo muchas herramientas modernas funcionan mediante capas de abstracción que desacoplan el hardware físico de la administración lógica, facilitando escalabilidad, automatización y tolerancia a fallos.
\subsection{Responsabilidades}
\label{sec:orga087be4}

\begin{itemize}
\item \textbf{Gael Emiliano:} Investigación técnica de herramientas utilizadas en laboratorio, documentación de comandos avanzados y explicación de procesos administrativos relacionados con almacenamiento, sincronización y servicios del sistema.
\item \textbf{Mariana:} Captura detallada de temas vistos en clase y ampliación conceptual relacionada con automatización, seguridad y administración de recursos.
\item \textbf{Isaac Julián:} Corrección técnica, integración final del archivo Org Mode, revisión de sintaxis y compilación del documento final.
\end{itemize}

\newpage
\section{Tema principal de la semana}
\label{sec:org5f7fd25}

\subsection{Conceptos nuevos}
\label{sec:orgcd7407e}

\subsubsection{Introducción general}
\label{sec:orge4dfee9}

Uno de los aspectos más importantes trabajados durante esta semana fue comprender cómo Linux administra servicios críticos mediante herramientas especializadas que interactúan directamente con el kernel, el sistema de archivos y los mecanismos internos de comunicación entre procesos.

A diferencia de otros sistemas operativos donde muchas tareas administrativas dependen de interfaces gráficas, Linux permite controlar prácticamente todos los componentes mediante comandos, archivos de configuración y servicios ejecutados en segundo plano.

Esta filosofía proporciona:

\begin{itemize}
\item Mayor control administrativo.
\item Automatización avanzada.
\item Configuraciones reproducibles.
\item Flexibilidad operativa.
\item Mejor capacidad de auditoría.
\item Escalabilidad administrativa.
\end{itemize}

Linux no busca ocultar el funcionamiento interno del sistema; al contrario, proporciona herramientas para administrarlo directamente.
\subsubsection{Funcionamiento interno}
\label{sec:orgbf02312}

Durante las prácticas se observó que muchas herramientas modernas funcionan mediante capas de abstracción.

Esto significa que:

\begin{itemize}
\item El usuario administra recursos lógicos.
\item Linux traduce operaciones hacia hardware físico.
\item El kernel coordina acceso, seguridad y sincronización.
\item Los servicios funcionan como intermediarios especializados.
\end{itemize}

Gracias a esto es posible:

\begin{itemize}
\item Redimensionar almacenamiento dinámicamente.
\item Compartir recursos de red.
\item Automatizar montajes.
\item Sincronizar relojes.
\item Delegar privilegios administrativos.
\item Monitorear procesos en tiempo real.
\end{itemize}

Todo esto puede realizarse sin detener completamente el sistema, lo cual representa una ventaja importante en entornos empresariales.
\subsubsection{Automatización y administración}
\label{sec:org4767934}

Uno de los conceptos más reforzados fue la automatización de tareas administrativas.
Linux utiliza:

\begin{itemize}
\item Demonios.
\item Servicios systemd.
\item Archivos de configuración.
\item Herramientas CLI.
\item Políticas de permisos.
\item Scripts administrativos.
\end{itemize}

Esto permite reducir intervención manual y mantener estabilidad operativa.

La automatización resulta especialmente importante en:

\begin{itemize}
\item Servidores empresariales.
\item Infraestructura de red.
\item Sistemas distribuidos.
\item Ambientes virtualizados.
\item Plataformas cloud.
\item Centros de datos.
\end{itemize}

La automatización no busca reemplazar al administrador, sino permitir configuraciones más consistentes y reproducibles.

\newpage
\section{Herramientas y comandos importantes}
\label{sec:orgc4e0c86}

\subsection{Administración del sistema}
\label{sec:org3624b25}

\begin{verbatim}
systemctl status
systemctl restart
systemctl enable
\end{verbatim}

Estos comandos permiten administrar servicios dentro del sistema Linux utilizando systemd.

Funciones principales:

\begin{itemize}
\item Consultar estado de servicios.
\item Reiniciar demonios.
\item Habilitar servicios al arranque.
\item Monitorear fallos.
\item Verificar disponibilidad de servicios.
\end{itemize}

\texttt{systemctl} se ha convertido en una de las herramientas centrales de administración moderna en Linux.
\subsection{Monitoreo y auditoría}
\label{sec:org7d6d556}

\begin{verbatim}
top
htop
journalctl
\end{verbatim}

Estas herramientas permiten monitorear:

\begin{itemize}
\item Consumo de CPU.
\item Uso de memoria RAM.
\item Logs del sistema.
\item Actividad de servicios.
\item Procesos problemáticos.
\item Estado general del sistema.
\end{itemize}

El monitoreo constante resulta fundamental para mantener estabilidad y detectar anomalías antes de que generen fallos mayores.
\subsubsection{top}
\label{sec:org73c4f34}

\texttt{top} permite visualizar procesos en tiempo real y observar consumo de recursos.

Es especialmente útil para:

\begin{itemize}
\item Detectar procesos con alto consumo.
\item Revisar carga del sistema.
\item Monitorear memoria.
\item Identificar cuellos de botella.
\end{itemize}
\subsubsection{htop}
\label{sec:orgde0f1c6}

\texttt{htop} es una versión más moderna e interactiva de \texttt{top}.

Ventajas principales:

\begin{itemize}
\item Interfaz más intuitiva.
\item Navegación sencilla.
\item Visualización más clara.
\item Mejor administración de procesos.
\end{itemize}
\subsubsection{journalctl}
\label{sec:org3dd6e82}

\texttt{journalctl} permite revisar logs generados por systemd.

Ejemplo:

\begin{verbatim}
journalctl -xe
\end{verbatim}

Esto facilita diagnóstico y auditoría del sistema.
\subsection{Administración de almacenamiento}
\label{sec:org91ce56f}

\begin{verbatim}
lsblk
df -h
mount
umount
\end{verbatim}

Estas herramientas permiten:

\begin{itemize}
\item Visualizar discos y particiones.
\item Verificar espacio disponible.
\item Montar sistemas de archivos.
\item Desmontar recursos activos.
\item Identificar dispositivos conectados.
\end{itemize}

Se reforzó la idea de que Linux representa dispositivos mediante archivos dentro de \texttt{/dev}.
\subsubsection{lsblk}
\label{sec:org3fb2a01}

\texttt{lsblk} muestra estructura de discos, particiones y dispositivos de almacenamiento.

\begin{verbatim}
lsblk
\end{verbatim}
\subsubsection{df -h}
\label{sec:org6b44b6a}

\texttt{df -h} permite revisar espacio utilizado y disponible en sistemas de archivos montados.

\begin{verbatim}
df -h
\end{verbatim}
\subsubsection{mount y umount}
\label{sec:orgce3e9c6}

\texttt{mount} conecta sistemas de archivos al árbol local de directorios.

\texttt{umount} permite desmontarlos correctamente.

Esto es fundamental para administración de recursos locales y remotos.
\subsection{Administración de permisos}
\label{sec:org2ffc371}

\begin{verbatim}
sudo
visudo
chmod
chown
\end{verbatim}

La administración correcta de permisos es uno de los pilares fundamentales de seguridad dentro de Linux.

Se enfatizó especialmente que:

\begin{itemize}
\item Cada privilegio otorgado tiene implicaciones de seguridad.
\item El principio de mínimo privilegio debe mantenerse.
\item La automatización no debe comprometer control administrativo.
\item Los permisos definen acceso a recursos críticos.
\end{itemize}
\subsubsection{chmod}
\label{sec:org5d0759a}

\texttt{chmod} permite modificar permisos de archivos y directorios.

\begin{verbatim}
chmod 755 archivo
\end{verbatim}
\subsubsection{chown}
\label{sec:org4cb3c76}

\texttt{chown} permite modificar propietario y grupo asociado a archivos.

\begin{verbatim}
chown usuario:grupo archivo
\end{verbatim}
\subsubsection{sudo y visudo}
\label{sec:orgac2c667}

\texttt{sudo} permite ejecutar tareas administrativas de manera controlada.

\texttt{visudo} permite modificar \texttt{/etc/sudoers} verificando sintaxis antes de guardar.

\newpage
\section{Conceptos de repaso}
\label{sec:orgc430f38}

\begin{itemize}
\item Linux utiliza arquitectura modular basada en herramientas especializadas.
\item El kernel funciona como intermediario entre hardware y software.
\item Los sistemas de archivos abstraen organización física de datos.
\item Los servicios funcionan mediante demonios ejecutados en segundo plano.
\item El almacenamiento moderno utiliza capas lógicas desacopladas del hardware físico.
\item Los permisos y usuarios son fundamentales para mantener estabilidad y seguridad.
\item La automatización permite reducir errores administrativos.
\end{itemize}
\section{Máximas}
\label{sec:orge55c9d9}

\begin{itemize}
\item “En Linux, automatizar no significa perder control; significa administrar el sistema con precisión.”
\item “Cada servicio ejecutándose en segundo plano representa una responsabilidad administrativa.”
\item “El kernel no improvisa: cada recurso del sistema pasa por políticas, permisos y planificación.”
\item “La diferencia entre un usuario común y un administrador no es ejecutar más comandos; es entender las consecuencias de ejecutarlos.”
\item “Linux abstrae el hardware para que el administrador piense en servicios, no en limitaciones físicas.”
\item “Todo sistema estable depende más de buenas configuraciones que de hardware poderoso.”
\end{itemize}
\section{Sección libre}
\label{sec:org523e7fa}

\begin{itemize}
\item Linux permite automatizar prácticamente cualquier tarea administrativa mediante scripts y servicios.
\item Los demonios trabajan constantemente administrando recursos críticos.
\item Los archivos dentro de \texttt{/etc} representan gran parte de la configuración persistente del sistema.
\item Muchas tecnologías modernas de Linux priorizan escalabilidad y tolerancia a fallos.
\item Las herramientas CLI ofrecen mayor precisión administrativa que muchas interfaces gráficas tradicionales.
\item El monitoreo continuo es esencial para mantener disponibilidad y estabilidad.
\end{itemize}

\newpage
\section{Realimentación y reflexión de aprendizajes}
\label{sec:orgf7bc584}

Durante esta semana se logró comprender con mayor profundidad cómo Linux organiza internamente la administración de recursos críticos y cómo múltiples tecnologías trabajan coordinadamente para mantener estabilidad, seguridad y automatización dentro del sistema operativo.

Uno de los aspectos más importantes observados fue que muchas herramientas aparentemente simples esconden gran complejidad interna relacionada con sincronización, permisos, abstracción de hardware y administración dinámica de recursos.

Esto permitió comprender que administrar Linux no consiste únicamente en memorizar comandos, sino en entender cómo interactúan entre sí los distintos componentes del sistema.

También se reforzó la importancia de la automatización y del monitoreo continuo dentro de ambientes profesionales. Servicios como systemd, herramientas de auditoría y configuraciones persistentes permiten mantener sistemas funcionando durante largos periodos sin intervención manual constante.

Finalmente, las prácticas ayudaron a desarrollar una visión más cercana a la administración real de servidores Linux, donde estabilidad, disponibilidad y seguridad dependen directamente de la calidad de las configuraciones realizadas por el administrador del sistema.

En conjunto, la semana permitió reforzar la idea de que Linux es un sistema operativo diseñado para administración precisa, automatización avanzada y escalabilidad en ambientes profesionales.
\end{document}
